\documentclass[10pt]{article}
\usepackage[T1]{fontenc}
\usepackage[utf8]{inputenc}
\usepackage[french]{babel}
\usepackage{newpxtext,newpxmath}
\usepackage[a4paper,top=1.5cm,bottom=1.35cm,left=1.6cm,right=1.6cm]{geometry}
\usepackage{microtype}
\usepackage{xcolor}
\usepackage{titlesec}
\usepackage[most]{tcolorbox}
\usepackage{enumitem}
\usepackage{array}
\usepackage{booktabs}
\usepackage{fancyhdr}

\definecolor{navy}{HTML}{1F3A5F}
\definecolor{grey}{HTML}{54606E}
\definecolor{green}{HTML}{2C6E49}
\definecolor{amber}{HTML}{9C6A16}
\definecolor{red}{HTML}{A6322E}
\definecolor{plum}{HTML}{7A4EA8}
\definecolor{boxgray}{HTML}{F3F5F8}
\definecolor{boxred}{HTML}{FCF2F1}
\definecolor{boxgreen}{HTML}{F1F7F3}
\definecolor{boxplum}{HTML}{F7F3FB}
\definecolor{boxcode}{HTML}{F7F8FA}
\definecolor{boxedge}{HTML}{D9E0EA}

\setlength{\parindent}{0pt}
\setlength{\parskip}{2.6pt}
\linespread{1.0}

\pagestyle{fancy}\fancyhf{}
\lhead{\small\color{grey}Scella}
\rhead{\small\color{grey}Les couples -- le croisement complet -- 11 août 2026}
\cfoot{\small\color{grey}\thepage}
\renewcommand{\headrulewidth}{0.3pt}

\titleformat{\section}{\color{navy}\large\bfseries}{}{0pt}{}
\titlespacing{\section}{0pt}{11pt}{3pt}
\titleformat{\subsection}{\color{navy}\bfseries}{}{0pt}{}
\titlespacing{\subsection}{0pt}{7pt}{1.5pt}

\newtcolorbox{these}{colback=boxgray,colframe=boxedge,boxrule=0.5pt,arc=2.5pt,
  left=8pt,right=8pt,top=4.5pt,bottom=4.5pt}
\newtcolorbox{alerte}{colback=boxred,colframe=red!55,boxrule=0.8pt,arc=2.5pt,
  left=8pt,right=8pt,top=4.5pt,bottom=4.5pt}
\newtcolorbox{acquis}{colback=boxgreen,colframe=green!45,boxrule=0.8pt,arc=2.5pt,
  left=8pt,right=8pt,top=4.5pt,bottom=4.5pt}
\newtcolorbox{piste}{colback=boxplum,colframe=plum!45,boxrule=0.8pt,arc=2.5pt,
  left=8pt,right=8pt,top=4.5pt,bottom=4.5pt}
\newtcolorbox{extrait}{colback=boxcode,colframe=boxedge,boxrule=0.5pt,arc=1.5pt,
  left=7pt,right=7pt,top=3.5pt,bottom=3.5pt}

\newcommand{\best}[1]{{\color{green}\textbf{#1}}}
\newcommand{\pire}[1]{{\color{red}\textbf{#1}}}
\newcommand{\sz}[1]{{\small\color{grey}\textsf{#1}}}
\newcommand{\code}[1]{\texttt{\small #1}}
\newcommand{\x}{$\times$}
\newcommand{\loc}[1]{{\footnotesize\color{grey}\texttt{#1}}}

\begin{document}

\begin{center}
{\Huge\bfseries\color{navy} Scella}\\[4pt]
{\large\color{grey}\itshape Les couples : le croisement complet de tous les éléments}\\[4pt]
{\small\color{grey} 11 août 2026 \; · \; cinquième et dernier document de la série « couples » \; · \; code lu à \code{aac2053}}
\end{center}
\vspace{2pt}

\begin{these}
\textbf{Le résultat central.} Les quatre documents précédents notent \textbf{huit couples choisis}. Le
croisement systématique des quatorze éléments -- 91 paires -- en fait apparaître \textbf{deux servis qui
n'étaient pas classés}, dont un qui prend la première place, et \textbf{un latent qui est le meilleur
rendement du dépôt}. Mais le vrai résultat est ailleurs, et il est inconfortable : le portefeuille
contient \textbf{neuf chaînes scellées} et \textbf{deux re-prouveurs à dépendances vides}. La propriété
qui fait toute la valeur de l'édifice -- \emph{un tiers re-prouve sans nous croire} -- couvre
aujourd'hui \textbf{deux chaînes sur neuf}.
\end{these}

\section{1 · Ce que la série couvre, et le trou exact}

Les quatre documents du jour partagent une méthode que je reprends sans la modifier : un couple est deux
éléments \emph{plus un liant}, c'est-à-dire le mécanisme qui leur fait produire ensemble une propriété
vérifiable qu'aucun n'a seul ; et la notation se fait sur les quatre lentilles intrinsèques du 5 août --
maturité, rareté, adversaire, démontrabilité -- depuis que la lentille d'objectif extérieur a disparu
avec la piste qu'elle servait.

Ce qui manque n'est donc pas la méthode, c'est la \textbf{couverture}. Les huit couples ont été choisis,
non obtenus par balayage. Sur quatorze éléments il existe 91 paires ; la série en documente huit, et rien
ne dit ce que valent les quatre-vingt-trois autres. Le présent document répond à cette question-là, et
seulement à celle-là : il ne remplace aucun des quatre, il les complète par le balayage qui manquait.

La plupart des paires balayées ne sont rien -- c'est attendu, et le §5 les nomme plutôt que de laisser
le silence passer pour de l'exhaustivité. Trois survivent, et une observation de portefeuille les domine
toutes.

\section{2 · L'inventaire réel de la surface scellée}

C'est le fait qui recadre l'ensemble, et aucun des quatre documents ne le mentionne.

\textbf{Neuf crates scellent une véritable chaîne HMAC-SHA256}, au sens fort : le vocabulaire de
chaînage y est présent (\code{chain}, \code{ledger}, \code{prev\_trace\_id}), pas seulement une primitive
cryptographique. Les voici :

\begin{extrait}
\code{efe-audit-core} \; · \; \code{efe-server} \; · \; \code{efe-gold-test-gate} \; · \;
\code{efe-publisher-identity} \; · \; \code{efe-nano} \\[1.5pt]
\code{efe-llm-router} \; · \; \code{efe-curation} \; · \; \code{efe-agent-runner} \; · \;
\code{efe-active-learning}
\end{extrait}

La recherche brute en renvoie onze. J'en ai écarté deux -- \code{efe-cache} et \code{efe-multi-region} --
après lecture : ils emploient HMAC à d'autres fins et ne chaînent rien. Annoncer onze aurait été
exactement la survente que ce projet s'interdit ; la différence entre « onze » et « neuf » est une
lecture de fichier, pas une nuance de présentation.

\renewcommand{\arraystretch}{1.15}
\begin{center}\small
\begin{tabular}{@{}>{\raggedright\arraybackslash}p{4.5cm} >{\raggedright\arraybackslash}p{4.3cm} c >{\raggedright\arraybackslash}p{5.6cm}@{}}
\toprule
\textbf{Chaîne scellée} & \textbf{Re-prouveur} & \textbf{0-dép} & \textbf{Ce que cela vaut} \\
\midrule
Bundle DAG (pile B) & \code{verify-mini} & \best{oui} & L'auditeur lit toute la source avant de croire \\
Reçu méta (Couche 3) & \code{verify\_meta\_mini} & \best{oui} & Idem, sur l'autre pile \\
Reçu méta (Couche 3) & \code{efe-meta-verify} & \pire{non} & 70 dépendances \\
\code{audit\_events} & \code{efe\_verify\_audit} & \pire{non} & Exige Postgres \emph{et} un fichier de clés \\
Registre gold & \code{efe\_gold\_verify} & \pire{non} & Tire tout le crate \\
Identité éditeur & \code{efe\_publisher\_verify} & \pire{non} & 115 dépendances \\
\textbf{Six autres chaînes} & \mbox{\pire{aucun}} & --- & Jamais re-prouvables par un tiers \\
\bottomrule
\end{tabular}
\end{center}

\begin{alerte}
\textbf{La conséquence, dite franchement.} Le couple classé n\textsuperscript{o}1 par la série --
« artefact scellé \x{} re-prouveur zéro-dépendance » -- a été \emph{correctement} identifié comme le
meilleur. Personne n'a mesuré sa \textbf{couverture} : il vaut pour \textbf{deux chaînes sur neuf}. Sept
chaînes scellées n'ont aucun re-prouveur qu'un tiers puisse lire intégralement avant de lui accorder sa
confiance -- et pour six d'entre elles, aucun re-prouveur du tout.
\end{alerte}

Il faut être juste : sceller sans re-prouveur tiers n'est pas inutile. Une chaîne HMAC reste
\emph{tamper-evident} pour le détenteur de la clé, et c'est déjà davantage que la plupart des systèmes.
Mais c'est précisément la propriété que le portefeuille met en avant, et l'écart entre « nous scellons
partout » et « un tiers peut re-prouver ici et ici » est le genre d'écart qu'un examinateur trouve seul.
Mieux vaut l'avoir écrit avant lui.

\section{3 · Les trois couples que les huit ont manqués}

\subsection{3.1 · C9 -- Pile A \x{} pile B, par la route bundle \sz{(servi)}}

Le rapport du 5 août désignait ce geste comme « le seul de tout le dépôt à transformer une addition en
multiplication », et le donnait à faire. \textbf{Il est fait.} \code{GET /v1/swarm/bundle/:id} exporte un
reçu de la pile A sous la forme exacte que \code{verify-mini} sait re-prouver -- et il porte trois
propriétés que la série ne relève nulle part, chacune étant une décision de conception et non un détail.

\textbf{L'attestation ed25519} de la racine n'est présente que si une clé d'attestation est déployée ;
sinon le champ est \textbf{absent}, jamais vide. La nuance décide de tout : un champ vide se confondrait
avec une signature, et un auditeur pressé conclurait à une attestation qui n'existe pas.

\textbf{La clé publique n'est jamais servie} à côté du bundle. La servir rendrait la vérification
circulaire -- n'importe qui peut engendrer une paire et signer ce qu'il veut. Seule l'empreinte
\code{key\_id} sort, pour que l'auditeur sache \emph{laquelle} de ses clés déjà épinglées appliquer.

\textbf{L'échec est fort.} Dérivation refusée $\Rightarrow$ 500 sans bundle partiel, parce qu'un bundle
amputé est indistinguable d'un bundle tronqué par un attaquant. Et si la publication de l'ancre rate, la
route rend une erreur plutôt que de servir un bundle annonçant \code{anchored: true} -- un auditeur
chercherait alors une ancre inexistante.

En toute honnêteté, ce n'est pas un couple entièrement neuf : c'est \textbf{l'instance concrète et
complète du couple n\textsuperscript{o}1}. La série a décrit le motif sans noter que le pont pile A
$\rightarrow$ pile B est le seul endroit du dépôt où il soit réalisé avec attestation \emph{et} ancre.
La distinction n'est pas verbale : c'est la différence entre un principe et une chose qui tourne.

\subsection{3.2 · C10 -- Un journal d'ancres, deux binaires indépendants \sz{(servi)}}

Celui-ci est réellement neuf. \code{efe-swarm-gateway} et \code{efe-meta-swarm-server} -- deux processus
distincts, deux piles de reçu sans code commun -- écrivent dans \textbf{le même} journal
\code{EFE\_SWARM\_ANCHOR\_LOG}. Le choix est délibéré et documenté dans le code lui-même
\loc{(controller/anchor.rs:31 : « le format et le fichier sont ceux du Fil P »)}.

La propriété composée qu'aucune moitié ne possède : un auditeur qui reçoit \textbf{un seul} journal peut
ordonner temporellement des événements provenant des \emph{deux} piles, et y détecter une troncature
dans l'une ou l'autre \textbf{sans détenir aucune clé}. Deux systèmes de scellement indépendants qui
partagent une ancre unique est une discipline peu courante -- l'inclination naturelle est de donner à
chaque sous-système son propre journal, ce qui interdit précisément l'ordonnancement croisé.

Sa limite est celle de tout l'ancrage aujourd'hui : le journal reste détenu par l'opérateur. Il supprime
la réécriture rétroactive, pas la confiance. C'est le jalon 3b, et il est moins coûteux qu'il n'y paraît
-- \code{forensic-export} porte déjà la structure \code{AnchorProof} et le contrôle de liaison de racine.

\subsection{3.3 · C11 -- Le registre gold \x{} un re-prouveur zéro-dép \sz{(latent)}}

\code{efe-gold-test-gate} scelle une chaîne HMAC avec \code{prev\_trace\_id}, et ce qu'elle scelle n'a
d'équivalent nulle part ailleurs dans le dépôt : \textbf{résultats de runs gold, usages de dérogation,
détections de dérive, et cérémonies d'approbation par un opérateur}.

La différence de nature mérite d'être posée. Les autres chaînes scellent des \emph{sorties} -- ce que le
système a produit. Celle-ci scelle \textbf{le processus de qualification lui-même} : qui a approuvé quel
jeu gold, qui a contourné la porte, et quand. C'est la première chose qu'un examinateur hostile cherche,
parce que c'est là que se logent les décisions humaines qu'on préfère oublier -- et c'est nettement plus
rare que de sceller des résultats.

\begin{piste}
\textbf{Il manque une seule pièce : un re-prouveur à dépendances vides.} \code{efe\_gold\_verify} existe
mais tire tout le crate, et il importe \code{InMemoryGoldRunLedger} -- point à contrôler avant d'en
dépendre. \textbf{Deux gabarits complets existent déjà} dans le dépôt (\code{verify-mini} et
\code{verify\_meta\_mini}, avec SHA-256, HMAC et parseur JSON écrits à la main). C'est le meilleur
rendement latent du portefeuille. \sz{Taille M.}
\end{piste}

Le corollaire corrige une erreur d'appréciation qui traîne depuis le 5 août. Le composant « banc de
mesure » y était noté \emph{« réfute sa propre thèse, rien n'est rejouable »}. C'est vrai de
\code{efe-bench}, qui ne scelle rien. Ce n'est \textbf{pas} vrai de \code{efe-gold-test-gate}, dont la
moitié difficile -- la chaîne scellée avec ses quatre natures d'événements -- est déjà faite. Le
composant avait été jugé comme un bloc alors qu'il en contient deux, de maturités opposées.

\section{4 · Le classement fusionné}

Les onze couples, notés sur les quatre lentilles intrinsèques, rangs de 1 à 11 par lentille.

\renewcommand{\arraystretch}{1.13}
\begin{center}\small
\begin{tabular}{@{}r >{\raggedright\arraybackslash}p{7.1cm} cccc c c@{}}
\toprule
& \textbf{Le couple} & \sz{Mat} & \sz{Rar} & \sz{Adv} & \sz{Dém} & \textbf{Moy} & \sz{Éc.} \\
\midrule
\textbf{1} & \textbf{Pile A \x{} pile B} \sz{(bundle + ed25519 + ancre)} & \best{1} & 2 & \best{1} & 3 & \textbf{1,75} & 2 \\
2 & Artefact scellé \x{} re-prouveur zéro-dép & 2 & 4 & 2 & 4 & \textbf{3,00} & 2 \\
3 & Topologie déclarée \x{} appariement injectif & 3 & 6 & 3 & 2 & \textbf{3,50} & 4 \\
\textbf{4} & \textbf{Ancre partagée, deux binaires} & 4 & 3 & 6 & 5 & \textbf{4,50} & 3 \\
5 & Liaison v4 \x{} ancre publiée & 5 & 7 & 5 & \best{1} & \textbf{4,50} & 6 \\
\textbf{6} & \textbf{Registre gold \x{} re-prouveur zéro-dép} & 7 & \best{1} & 4 & \pire{10} & \textbf{5,50} & \pire{9} \\
7 & Canari de Lipschitz \x{} porte de démarrage & 8 & 5 & 7 & 6 & 6,50 & 3 \\
8 & Caviardage PII \x{} scellement & 6 & 9 & 8 & 7 & 7,50 & 3 \\
9 & Plafonnement \x{} coût scellé & 9 & 10 & 9 & 8 & 9,00 & 2 \\
10 & Co-signature R3 \x{} chaîne HMAC & 10 & 8 & 10 & 9 & 9,25 & 2 \\
11 & Revendication \x{} sa propre chute & 11 & 11 & 11 & 11 & 11,00 & \best{0} \\
\bottomrule
\end{tabular}
\end{center}

\sz{En gras : les trois couples absents des quatre documents précédents. Éc. = amplitude entre le
meilleur et le pire rang.}

\begin{alerte}
\textbf{Faits et jugements, séparés.} Les mécanismes, emplacements et taxonomies de ce document sont des
\textbf{faits}, lus dans le code du 11 août ; l'annexe donne la commande qui rejoue chacun. Les
\textbf{notes} sont un \textbf{jugement} structuré pour être discutable. Un désaccord sur « rareté = 3 »
est légitime ; un désaccord sur « neuf chaînes scellées » s'ouvre en une commande.
\end{alerte}

Deux lectures comptent davantage que les rangs eux-mêmes.

\textbf{C11 a un écart de 9} -- premier en rareté, dernier en démontrabilité. C'est exactement le profil
\emph{pire réalisé, meilleur latent} que la série avait diagnostiqué sur la co-signature R3, mais dans
une version plus favorable : un seul fil le referme, là où R3 attend un producteur entier.

\textbf{C9 déloge l'ancien premier}, et pas parce que celui-ci était mal jugé. Le couple
n\textsuperscript{o}1 énonçait un motif ; C9 en est l'incarnation complète, avec deux propriétés que le
motif seul n'implique pas -- l'attestation non circulaire et l'ancre. Un principe et une chose qui
tourne ne se classent pas au même rang.

\section{5 · Les croisements morts, nommés}

Le balayage a surtout produit des non-couples. Les taire aurait laissé croire que l'exhaustivité était
gratuite.

\textbf{Fil B \x{} méta-essaim.} Aucun chemin partagé : ils vivent dans deux binaires distincts, à deux
couches différentes. C'est précisément l'objet du Fil C -- la Colonne -- et le Fil C fait \textbf{zéro
ligne}. La paire qu'on croit voir sur un schéma d'architecture n'existe pas dans le code.

\textbf{Runtime nano \x{} tout.} \code{efe-nano-inference-daemon} n'a aucun appelant dans le dépôt.
Toutes ses paires sont vides par construction, quel que soit le partenaire.

\textbf{Sandwich C.1, essaim médiateur, Fil C \x{} tout.} Zéro occurrence de \code{HopKind},
\code{SwarmMediator}, \code{EdgePlan}. Ce sont des conceptions, pas des moitiés de couple : on ne peut
pas apparier ce qui n'existe pas, et un document d'architecture ne compte pas comme une moitié.

\textbf{Flexible MLP \x{} tout.} Hors dépôt.

\begin{alerte}
\textbf{Une correction, au passage.} L'essaim $\Phi$ est classé dans la série comme si la pureté de sa
concurrence était acquise. \textbf{Elle ne l'est pas dans \code{HEAD}} : \code{dissolution\_centroid} y
compte \textbf{zéro occurrence}. Ces 323 lignes appartiennent à la session concurrente et vivent
uniquement dans l'arbre de travail. Tout couple qui s'appuie sur cette propriété est aujourd'hui
\emph{latent}, non \emph{servi}.
\end{alerte}

\section{6 · Le jugement}

Le portefeuille de couples est \textbf{plus riche et plus déséquilibré} que la série ne le montrait. Deux
couples servis n'étaient pas classés, dont celui qui prend la première place ; un latent de premier plan
dormait dans un composant jugé en bloc alors qu'il en contient deux.

Mais la conclusion qui compte n'est pas un rang. C'est que la propriété fondatrice de tout l'édifice --
\emph{un tiers re-prouve sans nous croire} -- \textbf{couvre deux chaînes scellées sur neuf}.

\begin{acquis}
\textbf{L'écart le plus rentable n'est plus de coder un couple nouveau.} C'est de \textbf{porter les deux
gabarits zéro-dépendance existants sur les chaînes qui n'en ont pas}, en commençant par le registre gold
-- la seule qui scelle un processus de décision \emph{humain} plutôt qu'un résultat machine. Le travail
est de la transcription, pas de l'invention : les deux gabarits sont écrits, testés, et leur arbre de
dépendances est vérifiablement vide.
\end{acquis}

\section{Annexe · Vérifier chaque affirmation soi-même}

Les affirmations \textbf{négatives} de ce document -- « six chaînes n'ont aucun re-prouveur », « le
Sandwich fait zéro ligne » -- sont les plus importantes à pouvoir contrôler, parce qu'un lecteur ne peut
pas les vérifier en lisant. Une recherche qui revient vide n'est une preuve qu'à condition d'être
rejouable. Les commandes ci-dessous ont toutes été exécutées avant publication, depuis la racine du
dépôt.

\renewcommand{\arraystretch}{1.22}
\begin{center}\small
\begin{tabular}{@{}>{\raggedright\arraybackslash}p{5.3cm} >{\raggedright\arraybackslash}p{10.8cm}@{}}
\toprule
\textbf{Affirmation} & \textbf{Ce qui la rejoue} \\
\midrule
Onze crates emploient HMAC \newline \sz{(neuf après lecture)} &
  \code{grep -rl 'Hmac<Sha256>\textbar HmacSha256' crates -{}-include='*.rs' \textbar\ cut -d/ -f2 \textbar\ sort -u} \\
Six binaires de vérification & \code{ls crates/*/src/bin/*.rs \textbar\ grep -i verif} \newline
  \sz{plus} \code{crates/verify-mini/src/main.rs} \sz{(nommé main, pas dans bin/)} \\
\code{verify-mini} est sans dépendance & \code{cargo tree -p verify-mini -e normal -{}-prefix none} \sz{→ une seule ligne} \\
\code{verify\_meta\_mini} aussi & \code{cargo tree -p verify-meta-mini -e normal -{}-prefix none} \\
Les autres vérifieurs sont lourds & \code{cargo tree -p efe-publisher-identity -e normal \textbar\ wc -l} \sz{→ 115} \\
C9 : la route bundle et ses gardes & \code{crates/efe-swarm-gateway/src/routes/bundle.rs} \sz{(en-tête, P.2 à P.4)} \\
C10 : l'ancre est partagée & \code{grep -rn 'EFE\_SWARM\_ANCHOR\_LOG\textbar ANCHOR\_LOG\_ENV' crates -{}-include='*.rs'} \\
C11 : ce que le registre gold scelle & \code{crates/efe-gold-test-gate/src/ledger/mod.rs} \sz{(en-tête, quatre natures)} \\
$\Phi$ pur non commité & \code{git show HEAD:crates/efe-core/src/els/mod.rs \textbar\ grep -c dissolution\_centroid} \sz{→ 0} \\
Sandwich / médiateur / Fil C & \code{grep -rn 'HopKind\textbar SwarmMediator\textbar EdgePlan' crates -{}-include='*.rs'} \sz{→ vide} \\
Runtime nano sans appelant & \code{grep -rn 'efe-nano-inference-daemon' crates/*/Cargo.toml} \\
\bottomrule
\end{tabular}
\end{center}

Deux avertissements pour que cette annexe ne soit pas lue plus largement qu'elle ne vaut. Les lignes
portant sur des \emph{absences} valent à la date du document ; et la classification « neuf chaînes sur
onze » repose sur une \textbf{lecture} des deux crates écartés, pas sur la commande -- c'est le seul
endroit du document où la commande seule ne suffit pas, et c'est pourquoi il est dit ici plutôt que
laissé à découvrir.

\vspace{4pt}
\begin{these}
\small\textbf{Rappel de cadrage.} Une chaîne scellée atteste l'\textbf{intégrité} d'une suite
d'enregistrements, jamais la \textbf{correction} de ce qu'ils décrivent. Elle est \emph{tamper-evident
pour un détenteur de clé}, pas \emph{trustless pour un tiers}. C'est vrai des neuf chaînes de ce
document, y compris des deux qui disposent d'un re-prouveur à dépendances vides : celui-ci change
\emph{qui} peut vérifier, pas \emph{ce} qui est attesté.
\end{these}

\end{document}
